\phantomsection
\addcontentsline{lot}{chapter}{\protect\numberline{}Resume Profil}

% Table generated by Excel2LaTeX from sheet 'Resume (unformatted)'
\begin{small}
\begin{longtable}{rY{12em}rrrrY{5em}Y{5em}}
    \multicolumn{8}{c}{RESUME PROFIL KESEHATAN}\\
    \multicolumn{8}{c}{KABUPATEN BELITUNG TIMUR}\\
    \multicolumn{8}{c}{TAHUN \tP}\\
	\\ \toprule

	\multirow{2}{*}{\textbf{No}} & \multirow{2}{*}{\textbf{INDIKATOR}}                & \multicolumn{5}{c}{\textbf{ANGKA/ NILAI}} & \multirow{2}{5em}{\raggedright\textbf{No. Lampiran}} \\
    \cmidrule{3-7}
	&                                                                                 & \textbf{L} & \textbf{P} & \textbf{L + P} & \textbf{Jumlah} & \textbf{Satuan} & \\ \midrule
	\endfirsthead
	\\ \toprule
	\multirow{2}{*}{\textbf{No}} & \multirow{2}{*}{\textbf{INDIKATOR}}                & \multicolumn{5}{c}{\textbf{ANGKA/ NILAI}} & \multirow{2}{5em}{\raggedright\textbf{No. Lampiran}} \\
    \cmidrule{3-7}
    &                                                                                 & \textbf{L} & \textbf{P} & \textbf{L + P} & \textbf{Jumlah} & \textbf{Satuan} & \\ \midrule
    \endhead
    % make bottom border on when table continue to next page
    \midrule
    \endfoot
    \endlastfoot
	           \textbf{I} & \textbf{GAMBARAN UMUM}                                                                                  &        &        &         &                    &                                &          \\
	                    1 & Luas Wilayah                                                                                            &        &        &         &              2.507 & Km\textsuperscript{2}          & Tabel 1  \\
	  \rowcolor{black!5}2 & Jumlah Desa/ Kelurahan                                                                                  &        &        &         &                 39 & Desa/Kelurahan                 & Tabel 1  \\
	                    3 & Jumlah Penduduk                                                                                         & 69.252 & 65.573 & 134.825 &                    & Jiwa                           & Tabel 2  \\
	  \rowcolor{black!5}4 & Rata-rata jiwa/ rumah tangga                                                                            &        &        &         &               3,23 & Jiwa                           & Tabel 1  \\
	                    5 & Kepadatan Penduduk /Km\textsuperscript{2}                                                               &        &        &         &              53,78 & Jiwa/Km\textsuperscript{2}     & Tabel 1  \\
	  \rowcolor{black!5}6 & Rasio Beban Tanggungan                                                                                  &        &        &         &              45,70 & per 100 penduduk produktif     & Tabel 2  \\
	                    7 & Rasio Jenis Kelamin                                                                                     &        &        &         &             105,61 &                                & Tabel 2  \\
	  \rowcolor{black!5}8 & Persentase Penduduk 15 tahun ke atas melek huruf                                                        &   0,00 &   0,00 &   99,14 &                    & \%                             & Tabel 3  \\
	                      &                                                                                                         &        &        &         &                    &                                &          \\
	          \textbf{II} & \textbf{SARANA KESEHATAN}                                                                               &        &        &         &                    &                                &          \\
	        \textbf{II.1} & \textbf{Sarana Kesehatan}                                                                               &        &        &         &                    &                                &          \\
	                    9 & Jumlah Rumah Sakit Umum                                                                                 &        &        &         &                  1 & RS                             & Tabel 4  \\
	 \rowcolor{black!5}10 & Jumlah Rumah Sakit Khusus                                                                               &        &        &         &                  0 & RS                             & Tabel 4  \\
	                   11 & Jumlah Puskesmas Rawat Inap                                                                             &        &        &         &                  4 & Puskesmas                      & Tabel 4  \\
	 \rowcolor{black!5}12 & Jumlah Puskesmas non-Rawat Inap                                                                         &        &        &         &                  3 & Puskesmas                      & Tabel 4  \\
	                   13 & Jumlah Puskesmas Keliling                                                                               &        &        &         &                  0 & Puskesmas keliling             & Tabel 4  \\
	 \rowcolor{black!5}14 & Jumlah Puskesmas pembantu                                                                               &        &        &         &                 39 & Pustu                          & Tabel 4  \\
	                   15 & Jumlah Apotek                                                                                           &        &        &         &                 29 & Apotek                         & Tabel 4  \\
	 \rowcolor{black!5}16 & Jumlah Klinik Pratama                                                                                   &        &        &         &                  6 & Klinik Pratama                 & Tabel 4  \\
	                   17 & Jumlah Klinik Utama                                                                                     &        &        &         &                  1 & Klinik Utama                   & Tabel 4  \\
	                      &                                                                                                         &        &        &         &                    &                                &          \\
	        \textbf{II.2} & \multicolumn{3}{Y{21em}}{\textbf{Akses dan Mutu Pelayanan Kesehatan}}                                                     &         &                    &                                &          \\
	                   18 & Cakupan Kunjungan Rawat Jalan                                                                           & 148,42 & 235,90 &  190,97 &                    & \%                             & Tabel 5  \\
	 \rowcolor{black!5}19 & Cakupan Kunjungan Rawat Inap                                                                            &   7,86 &  11,99 &    9,87 &                    & \%                             & Tabel 5  \\
	                   20 & Angka kematian kasar/ \textit{Gross Death Rate} (GDR) di RS                                             &  61,26 &  39,20 &   48,39 &                    & per 1.000 pasien keluar        & Tabel 6  \\
	 \rowcolor{black!5}21 & Angka kematian murni/ \textit{Nett Death Rate} (NDR) di RS                                              &  28,02 &  17,87 &   22,09 &                    & per 1.000 pasien keluar        & Tabel 6  \\
	                   22 & Persentase \textit{Bed Occupation Rate} (BOR) di RS                                                     &        &        &         &              54,41 & \%                             & Tabel 7  \\
	 \rowcolor{black!5}23 & \textit{Bed Turn Over (BTO)} di RS                                                                      &        &        &         &              58,43 & Kali                           & Tabel 7  \\
	                   24 & \textit{Turn of Interval (TOI)} di RS                                                                   &        &        &         &               2,85 & Hari                           & Tabel 7  \\
	 \rowcolor{black!5}25 & \textit{Average Length of Stay (ALOS)} di RS                                                            &        &        &         &               3,40 & Hari                           & Tabel 7  \\
	                   26 & Persentase Puskesmas dengan ketersediaan obat esensial \& vaksin IRL                                    &        &        &         &             100,00 & \%                             & Tabel 11 \\
	                      &                                                                                                         &        &        &         &                    &                                &          \\
	        \textbf{II.3} & \multicolumn{3}{Y{21em}}{\textbf{Upaya Kesehatan Bersumberdaya Masyarakat (UKBM)}}                                        &         &                    &                                &          \\
	                   27 & Jumlah Posyandu Siklus Hidup                                                                            &        &        &         &                134 & Posyandu                       & Tabel 12 \\
	 \rowcolor{black!5}28 & Persentase Posyandu Siklus Hidup Aktif                                                                  &        &        &         &             100,00 & \%                             & Tabel 12 \\
	                      &                                                                                                         &        &        &         &                    &                                &          \\
	         \textbf{III} & \multicolumn{3}{Y{21em}}{\textbf{SUMBER DAYA MANUSIA KESEHATAN}}                                                          &         &                    &                                &          \\
	                   29 & Rasio Dokter Spesialis per 1.000 Penduduk                                                               &        &        &         &              34,12 & per 1.000 penduduk             & Tabel 13 \\
	 \rowcolor{black!5}30 & Rasio Dokter Sub Spesialis per 1.000 Penduduk                                                           &        &        &         &               7,42 & per 1.000 penduduk             & Tabel 13 \\
	                   31 & Rasio Dokter per 1.000 Penduduk                                                                         &        &        &         &               0,74 & per 1.000 penduduk             & Tabel 13 \\
	 \rowcolor{black!5}32 & Rasio Dokter Gigi Spesialis per 1.000 Penduduk                                                          &        &        &         &               5,19 & per 1.000 penduduk             & Tabel 13 \\
	                   33 & Rasio Dokter Gigi Sub Spesialis per 1.000 Penduduk                                                      &        &        &         &               0,74 & per 1.000 penduduk             & Tabel 13 \\
	 \rowcolor{black!5}34 & Rasio Dokter Gigi per 1.000 Penduduk                                                                    &        &        &         &               0,00 & per 1.000 penduduk             & Tabel 13 \\
	                   35 & Rasio Keperawatan per 1.000 Penduduk                                                                    &        &        &         &             278,88 & per 1.000 penduduk             & Tabel 14 \\
	 \rowcolor{black!5}36 & Rasio Tenaga Kebidanan per 1.000 Penduduk                                                               &        &        &         &             114,22 & per 1.000 penduduk             & Tabel 14 \\
	                   37 & Rasio Tenaga Kesehatan Masyarakat per 1.000 Penduduk                                                    &        &        &         &              21,51 & per 1.000 penduduk             & Tabel 15 \\
	 \rowcolor{black!5}38 & Rasio Tenaga Kesehatan Lingkungan per 1.000 Penduduk                                                    &        &        &         &               6,68 & per 1.000 penduduk             & Tabel 15 \\
	                   39 & Rasio Tenaga Gizi per 1.000 Penduduk                                                                    &        &        &         &              16,32 & per 1.000 penduduk             & Tabel 15 \\
	 \rowcolor{black!5}40 & Rasio Tenaga Kefarmasian per 1.000 Penduduk                                                             &        &        &         &              27,44 & per 1.000 penduduk             & Tabel 16 \\
	                   41 & Rasio Tenaga Psikologis Klinis per 1.000 Penduduk                                                       &        &        &         &               0,74 & per 1.000 penduduk             & Tabel 16 \\
	 \rowcolor{black!5}42 & Rasio Tenaga Kesehatan Tradisional per 1.000 Penduduk                                                   &        &        &         &               0,00 & per 1.000 penduduk             & Tabel 16 \\
	                   43 & Rasio Tenaga Tehnik Biomedika per 1.000 Penduduk                                                        &        &        &         &              20,77 & per 1.000 penduduk             & Tabel 17 \\
	 \rowcolor{black!5}44 & Rasio Tenaga Tehnik Keterapian Fisik per 1.000 Penduduk                                                 &        &        &         &               8,16 & per 1.000 penduduk             & Tabel 17 \\
	                   45 & Rasio Tenaga Keteknisan Medis per 1.000 Penduduk                                                        &        &        &         &              25,96 & per 1.000 penduduk             & Tabel 17 \\
	                      &                                                                                                         &        &        &         &                    &                                &          \\
	          \textbf{IV} & \multicolumn{3}{Y{21em}}{\textbf{PEMBIAYAAN KESEHATAN}}                                                                   &         &                    &                                &          \\
	                   46 & Cakupan Jaminan Kesehatan Nasional (JKN)                                                                &        &        &         &              98,45 & \%                             & Tabel 19 \\
	 \rowcolor{black!5}47 & Total anggaran kesehatan                                                                                &        \multicolumn{4}{r}{ 220.555.334.822,00} & Rp                             & Tabel 20 \\
	                   48 & Persentase APBD kesehatan terhadap APBD Kab.                                                            &        &        &         &              25,94 & \%                             & Tabel 20 \\
	 \rowcolor{black!5}49 & Anggaran kesehatan perkapita                                                                            &               \multicolumn{4}{r}{1.635.863,79} & Rp                             & Tabel 20 \\
	                      &                                                                                                         &        &        &         &                    &                                &          \\
						  &                                                                                                         &        &        &         &                    &                                &          \\
	           \textbf{V} & \textbf{KESEHATAN KELUARGA}                                                                             &        &        &         &                    &                                &          \\
	         \textbf{V.1} & \textbf{Kesehatan Ibu}                                                                                  &        &        &         &                    &                                &          \\
	                   50 & Jumlah Lahir Hidup                                                                                      &    877 &    786 &   1.663 &                    & Orang                          & Tabel 21 \\
	 \rowcolor{black!5}51 & Angka Lahir Mati (dilaporkan)                                                                           &        &        &   10,12 &                    & per 1.000 Kelahiran Hidup      & Tabel 21 \\
	                   52 & Jumlah Kematian Ibu                                                                                     &        &      4 &         &                    & Ibu                            & Tabel 23 \\
	 \rowcolor{black!5}53 & Angka Kematian Ibu (dilaporkan)                                                                         &        & 240,53 &         &                    & per 100.000 Kelahiran Hidup    & Tabel 24 \\
	                   54 & Persentase Kunjungan Ibu Hamil (K1)                                                                     &        &  97,18 &         &                    & \%                             & Tabel 26 \\
	 \rowcolor{black!5}55 & Persentase Kunjungan Ibu Hamil (K6)                                                                     &        &  90,50 &         &                    & \%                             & Tabel 26 \\
	                   56 & Persentase Persalinan di Fasyankes                                                                      &        &  90,34 &         &                    & \%                             & Tabel 26 \\
	 \rowcolor{black!5}57 & Persentase Pelayanan Ibu Nifas KF Lengkap                                                               &        &  90,50 &         &                    & \%                             & Tabel 26 \\
	                   58 & Persentase Ibu Nifas Mendapat Vitamin A                                                                 &        &  90,50 &         &                    & \%                             & Tabel 26 \\
	 \rowcolor{black!5}59 & Persentase Ibu hamil dengan imunisasi Td2+                                                              &        &  48,21 &         &                    & \%                             & Tabel 27 \\
	                   60 & Persentase Ibu Hamil Mengonsumsi Tablet Tambah Darah 180 Tablet                                         &        &  84,09 &         &                    & \%                             & Tabel 28 \\
	 \rowcolor{black!5}61 & Persentase Bumil dengan Komplikasi Kebidanan yang Ditangani                                             &        & 170,20 &         &                    & \%                             & Tabel 32 \\
	                   62 & Persentase Peserta KB Aktif Modern                                                                      &        &        &   85,61 &                    & \%                             & Tabel 29 \\
	 \rowcolor{black!5}63 & Persentase Peserta KB Pasca Persalinan                                                                  &        &        &   76,06 &                    & \%                             & Tabel 31 \\
	                      &                                                                                                         &        &        &         &                    &                                &          \\
	         \textbf{V.2} & \textbf{Kesehatan Anak}                                                                                 &        &        &         &                    &                                &          \\
	                   64 & Jumlah Kematian Neonatal                                                                                &      3 &      8 &      11 &                    & neonatal                       & Tabel 34 \\
	 \rowcolor{black!5}65 & Angka Kematian Neonatal (dilaporkan)                                                                    &        &        &    6,61 &                    & per 1.000 Kelahiran Hidup      & Tabel 34 \\
	                   66 & Jumlah Bayi Mati                                                                                        &      7 &     11 &      18 &                    & bayi                           & Tabel 34 \\
	 \rowcolor{black!5}67 & Angka Kematian Bayi (dilaporkan)                                                                        &        &        &   10,82 &                    & per 1.000 Kelahiran Hidup      & Tabel 34 \\
	                   68 & Jumlah Balita Mati                                                                                      &      2 &      1 &       3 &                    & Balita                         & Tabel 34 \\
	 \rowcolor{black!5}69 & Persentase Angka Kematian Balita (dilaporkan)                                                           &        &        &    1,80 &                    & per 1.000 Kelahiran Hidup      & Tabel 34 \\
	                   70 & Persentase Bayi Baru Lahir Ditimbang                                                                    & 100,00 & 100,00 &  100,00 &                    & \%                             & Tabel 38 \\
	 \rowcolor{black!5}71 & Persentase Berat Badan Bayi Lahir Rendah (BBLR)                                                         &   9,12 &  11,32 &   10,16 &                    & \%                             & Tabel 38 \\
	                   72 & Persentase Kunjungan Neonatus 1 (KN 1)                                                                  &  99,89 &  99,62 &   99,76 &                    & \%                             & Tabel 39 \\
	 \rowcolor{black!5}73 & Persentase Kunjungan Neonatus 3 kali (KN Lengkap)                                                       &  98,75 &  99,11 &   98,92 &                    & \%                             & Tabel 39 \\
	                   74 & Persentase Bayi yang diberi ASI Eksklusif                                                               &        &        &   75,20 &                    & \%                             & Tabel 40 \\
	 \rowcolor{black!5}75 & Cakupan Imunisasi Campak/ Rubela pada Bayi                                                              &  86,52 &  87,15 &   86,82 &                    & \%                             & Tabel 43 \\
	                   76 & Persentase Imunisasi dasar lengkap pada bayi                                                            &  86,73 &  88,66 &   87,64 &                    & \%                             & Tabel 43 \\
	 \rowcolor{black!5}77 & Cakupan Bayi Mendapat Vitamin A                                                                         &        &        &   99,08 &                    & \%                             & Tabel 46 \\
	                   78 & Cakupan Anak Balita Mendapat Vitamin A                                                                  &        &        &   99,05 &                    & \%                             & Tabel 46 \\
	 \rowcolor{black!5}79 & Persentase Balita Mendapatkan Vitamin A                                                                 &        &        &   99,08 &                    & \%                             & Tabel 46 \\
	                   80 & Persentase Balita Memiliki Buku KIA                                                                     &        &        &   89,42 &                    & \%                             & Tabel 47 \\
	 \rowcolor{black!5}81 & Persentase Balita Dipantau Pertumbuhan dan Perkembangan                                                 &        &        &   90,97 &                    & \%                             & Tabel 47 \\
	                   82 & Persentase Balita Ditimbang (D/S)                                                                       &  76,75 &  79,63 &   78,18 &                    & \%                             & Tabel 48 \\
	 \rowcolor{black!5}83 & Persentase Balita Gizi Kurang (BB/TB)                                                                   &        &        &    2,96 &                    & \%                             & Tabel 49 \\
	                   84 & Persentase Balita Gizi Buruk (BB/TB)                                                                    &        &        &    0,01 &                    & \%                             & Tabel 49 \\
	 \rowcolor{black!5}85 & Cakupan Pemeriksaan Kesehatan Gratis Siswa Kelas 1 SD/ MI                                               &        &        &  100,00 &                    & \%                             & Tabel 50 \\
	                   86 & Cakupan Pemeriksaan Kesehatan Gratis Siswa Kelas 7 SMP/ MTs                                             &        &        &  100,00 &                    & \%                             & Tabel 50 \\
	 \rowcolor{black!5}87 & Cakupan Pemeriksaan Kesehatan Gratis Siswa Kelas 10 SMA/ MA                                             &        &        &  100,00 &                    & \%                             & Tabel 50 \\
	                   88 & Cakupan Pemeriksaan Kesehatan Gratis pada usia pendidikan dasar                                         &        &        &  100,00 &                    & \%                             & Tabel 50 \\
	 \rowcolor{black!5}89 & Cakupan Imunisasi HPV                                                                                   &        &        &   93,23 &                    & \%                             & Tabel 51 \\
	                   90 & Cakupan Imunisasi Anak Sekolah Lengkap                                                                  &  92,72 &  93,77 &   93,23 &                    & \%                             & Tabel 51 \\
	 \rowcolor{black!5}91 & Cakupan Pelayanan Kesehatan Gigi dan Mulut                                                              &        &        &    7,05 &                    & \%                             & Tabel 52 \\
	                   92 & Cakupan Pelayanan Kesehatan Gigi dan Mulut Pada Anak SD atau Setara                                     &        &        &   45,91 &                    & \%                             & Tabel 53 \\
	                      &                                                                                                         &        &        &         &                    &                                &          \\
	         \textbf{V.3} & \multicolumn{3}{Y{21em}}{\textbf{Kesehatan Usia Produktif dan Usia Lanjut}}                                               &         &                    &                                &          \\
	                   93 & Persentase Pelayanan Kesehatan Usia Produktif                                                           &        &        &   96,56 &                    & \%                             & Tabel 54 \\
	 \rowcolor{black!5}94 & Persentase Catin Mendapatkan Layanan Kesehatan                                                          &        &        &  100,00 &                    & \%                             & Tabel 55 \\
	                   95 & Persentase Pelayanan Kesehatan Usila (60+ tahun)                                                        &        &        &   82,42 &                    & \%                             & Tabel 56 \\
	                      &                                                                                                         &        &        &         &                    &                                &          \\
	          \textbf{VI} & \multicolumn{3}{Y{21em}}{\textbf{PENGENDALIAN PENYAKIT}}                                                                  &         &                    &                                &          \\
	        \textbf{VI.1} & \multicolumn{3}{Y{21em}}{\textbf{Pengendalian Penyakit Menular Langsung}}                                                 &         &                    &                                &          \\
	                   96 & Persentase orang terduga TBC mendapatkan pelayanan kesehatan sesuai standar                             &        &        &  133,71 &                    & \%                             & Tabel 59 \\
	 \rowcolor{black!5}97 & Cakupan Penemuan Kasus Tuberkulosis  (\%)                                                               &        &        &   90,93 &                    & \%                             & Tabel 59 \\
	                   98 & Cakupan Pemberian Terapi Pencegahan TB Pada Kontak Serumah                                              &        &        &  381,00 &                    & \%                             & Tabel 59 \\
	 \rowcolor{black!5}99 & Persentase angka pengobatan lengkap semua kasus TBC                                                     &  32,17 &  32,19 &   52,66 &                    & \%                             & Tabel 60 \\
	                  100 & Persentase angka keberhasilan pengobatan (\textit{Success Rate}) semua kasus TBC                        &  53,91 &  50,68 &   52,66 &                    & \%                             & Tabel 60 \\
	\rowcolor{black!5}101 & Persentase jumlah kematian selama pengobatan tuberkulosis                                               &        &        &    8,78 &                    & \%                             & Tabel 60 \\
	                  102 & Persentase penemuan penderita pneumonia pada balita                                                     &        &        &   18,60 &                    & \%                             & Tabel 61 \\
	\rowcolor{black!5}103 & Persentase puskesmas yang melakukan tatalaksana standar pneumonia min 60\%                              &        &        &         &             100,00 & \%                             & Tabel 61 \\
	                  104 & Jumlah Kasus HIV                                                                                        &      9 &      6 &      15 &                    & Kasus                          & Tabel 62 \\
	\rowcolor{black!5}105 & Persentase ODHIV Baru Mendapat Pengobatan ARV                                                           &        &        &   93,33 &                    & \%                             & Tabel 63 \\
	                  106 & Persentase Penderita Diare pada Semua Umur Dilayani                                                     &        &        &   54,20 &                    & \%                             & Tabel 64 \\
	\rowcolor{black!5}107 & Persentase Penderita Diare pada Balita Dilayani                                                         &        &        &   27,25 &                    & \%                             & Tabel 64 \\
	                  108 & Persentase Ibu hamil diperiksa Hepatitis                                                                &        &        &   91,31 &                    & \%                             & Tabel 65 \\
	\rowcolor{black!5}109 & Persentase Ibu hamil diperiksa Reaktif Hepatitis                                                        &        &        &    2,02 &                    & \%                             & Tabel 65 \\
	                  110 & Persentase Bayi dari Bumil Reakif Hepatitis Diperiksa                                                   &        &        &  100,00 &                    & \%                             & Tabel 66 \\
	\rowcolor{black!5}111 & Jumlah Kasus Baru Kusta (PB+MB)                                                                         &      4 &      2 &       6 &                    & Kasus                          & Tabel 67 \\
	                  112 & Angka penemuan kasus baru kusta (NCDR)                                                                  &   5,78 &   3,05 &    4,45 &                    & per 100.000 penduduk           & Tabel 67 \\
	\rowcolor{black!5}113 & Persentase Kasus Baru Kusta anak < 15 Tahun                                                             &        &        &   50,00 &                    & \%                             & Tabel 68 \\
	                  114 & Persentase Cacat Tingkat 0 Penderita Kusta                                                              &        &        &   50,00 &                    & \%                             & Tabel 68 \\
	\rowcolor{black!5}115 & Persentase Cacat Tingkat 2 Penderita Kusta                                                              &        &        &    0,00 &                    & \%                             & Tabel 68 \\
	                  116 & Angka Cacat Tingkat 2 Penderita Kusta                                                                   &        &        &   22,25 &                    & per 1.000.000 penduduk         & Tabel 68 \\
	\rowcolor{black!5}117 & Angka Prevalensi Kusta                                                                                  &        &        &    0,45 &                    & per 10.000 Penduduk            & Tabel 69 \\
	                  118 & Persentase Penderita Kusta PB Selesai Berobat (RFT PB)                                                  &        &        &    NULL &                    & \%                             & Tabel 70 \\
	\rowcolor{black!5}119 & Persentase Penderita Kusta MB Selesai Berobat (RFT MB)                                                  &        &        &  100,00 &                    & \%                             & Tabel 70 \\
	                      &                                                                                                         &        &        &         &                    &                                &          \\
	        \textbf{VI.2} & \multicolumn{3}{Y{21em}}{\textbf{Pengendalian Penyakit yang Dapat Dicegah dengan Imunisasi}}                              &         &                    &                                &          \\
	                  120 & AFP Rate (non polio) < 15 tahun                                                                         &        &        &    0,00 &                    & per 100.000 penduduk <15 tahun & Tabel 71 \\
	\rowcolor{black!5}121 & Jumlah kasus difteri                                                                                    &      0 &      0 &       0 &                    & Kasus                          & Tabel 72 \\
	                  122 & Persentase \textit{Case Fatality Rate} difteri                                                          &        &        &    NULL &                    & \%                             & Tabel 72 \\
	\rowcolor{black!5}123 & Jumlah kasus pertusis                                                                                   &      0 &      0 &       0 &                    & Kasus                          & Tabel 72 \\
	                  124 & Jumlah kasus tetanus neonatorum                                                                         &      0 &      0 &       0 &                    & Kasus                          & Tabel 72 \\
	\rowcolor{black!5}125 & Persentase \textit{Case Fatality Rate} tetanus neonatorum                                               &        &        &    NULL &                    & \%                             & Tabel 72 \\
	                  126 & Jumlah kasus hepatitis B                                                                                &      0 &      0 &       0 &                    & Kasus                          & Tabel 72 \\
	\rowcolor{black!5}127 & Jumlah kasus suspek campak                                                                              &      0 &      0 &       0 &                    & Kasus                          & Tabel 72 \\
	                  128 & Insiden rate suspek campak                                                                              &   0,00 &   0,00 &    0,00 &                    & per 100.000 penduduk           & Tabel 72 \\
	\rowcolor{black!5}129 & Persentase KLB ditangani < 24 jam                                                                       &        &        &         &              100,0 & \%                             & Tabel 73 \\
	                      &                                                                                                         &        &        &         &                    &                                &          \\
	        \textbf{VI.3} & \multicolumn{3}{Y{21em}}{\textbf{Pengendalian Penyakit Tular Vektor dan Zoonotik}}                                        &         &                    &                                &          \\
	                  130 & Angka kesakitan (Incidence Rate)DBD                                                                     &        &        &   35,60 &                    & per 100.000 penduduk           & Tabel 75 \\
	\rowcolor{black!5}131 & Persentase Angka kematian (\textit{Case Fatality Rate}) DBD                                             &   0,00 &   0,00 &    0,00 &                    & \%                             & Tabel 75 \\
	                  132 & Angka kesakitan malaria (\textit{Annual Parasite Incidence})                                            &        &        &    0,01 &                    & per 1.000 penduduk             & Tabel 76 \\
	\rowcolor{black!5}133 & Persentase konfirmasi laboratorium pada suspek malaria                                                  &        &        &       1 &                    & \%                             & Tabel 76 \\
	                  134 & Persentase pengobatan standar kasus malaria positif                                                     &        &        &       1 &                    & \%                             & Tabel 76 \\
	\rowcolor{black!5}135 & Penderita kronis filariasis                                                                             &      7 &      1 &       8 &                    & Kasus                          & Tabel 77 \\
	                      &                                                                                                         &        &        &         &                    &                                &          \\
	        \textbf{VI.4} & \multicolumn{3}{Y{21em}}{\textbf{Pengendalian Penyakit Tidak Menular}}                                                    &         &                    &                                &          \\
	                  136 & Persentase penderita Hipertensi Mendapat Pelayanan Kesehatan                                            &  50,12 & 121,69 &   85,02 &                    & \%                             & Tabel 78 \\
	\rowcolor{black!5}137 & Persentase penyandang DM yang terkendali                                                                &        &        &   51,15 &                    & \%                             & Tabel 79 \\
	                  138 & Persentase HPV (+) yang dilakukan IVA pada perempuan usia 30-69 tahun                                   &        &   NULL &         &                    & \% perempuan usia 30-69 tahun  & Tabel 80 \\
	\rowcolor{black!5}139 & Persentase HPV (+) dan IVA (+) yang dilakukan tata laksana                                              &        &   NULL &         &                    & \%                             & Tabel 80 \\
	                  140 & Persentase skrining payudara dengan SADANIS pada perempuan 30-69 tahun                                  &        &  15,00 &         &                    & \%                             & Tabel 80 \\
	\rowcolor{black!5}141 & Persentase skrining payudara dengan USG payudara pada perempuan 30-69 tahun                             &        &   0,00 &         &                    & \%                             & Tabel 80 \\
	                  142 & Persentase pelayanan Kesehatan Orang dengan Gangguan Jiwa Berat                                         &        &        &  100,00 &                    & \%                             & Tabel 81 \\
	                      &                                                                                                         &        &        &         &                    &                                &          \\
	         \textbf{VII} & \multicolumn{3}{Y{21em}}{\textbf{KESEHATAN LINGKUNGAN}}                                                                   &         &                    &                                &          \\
	                  143 & Persentase sarana air minum yang memenuhi syarat/ Diperiksa Kualitas Air Minumnya Sesuai Standar (Aman) &        &        &         &              63,64 & \%                             & Tabel 82 \\
	\rowcolor{black!5}144 & Persentase rumah tangga dengan air minum yang memenuhi syarat                                           &        &        &         &              55,24 & \%                             & Tabel 83 \\
	                  145 & Persentase KK dengan Akses terhadap Fasilitas Sanitasi                                                  &        &        &         &             100,00 & \%                             & Tabel 84 \\
	\rowcolor{black!5}146 & Persentase KK Stop BABS (SBS)                                                                           &        &        &         &             100,00 & \%                             & Tabel 85 \\
	                  147 & Persentase Desa/ Kelurahan 5 Pilar STBM                                                                 &        &        &         &               7,69 & \%                             & Tabel 85 \\
	\rowcolor{black!5}148 & Persentase Tempat Fasilitas Umum (TFU)  yang dilakukan pengawasan sesuai standar                        &        &        &         &              94,23 & \%                             & Tabel 88 \\
	                  149 & Persentase Tempat Pengelolaan Pangan (TPP) yang memenuhi syarat kesehatan                               &        &        &         &              55,81 & \%                             & Tabel 87 \\
	\rowcolor{black!5}150 & Persentase Hasil Pengukuran Kualitas Udara dalam Ruang Memenuhi Syarat                                  &        &        &         &               0,00 & \%                             & Tabel 88
\end{longtable}%
\end{small}
